\section{A 239-state machine by splitting an entry context}
\label{sec:fold239}

The decisive reduction is a change of control representation in a small
region of the repaired boundary240 table.  State $218$ has two meanings
distinguished by the bit immediately left of the head.  Splitting these
entry contexts allows one meaning to share state $179$, while the other
uses state $223$, which previously served only as an internal projection
helper.  Thus three old control states are represented by two new states.
The representation change needs no extra tape symbol, no extra tape cell,
and no assumption about the register values or the source program counter.

\subsection{Exact construction}

Let $M$ be the repaired boundary240 machine, with inherited start label
$95$ and halt label $254$.  Its SHA-256 digest is
\begin{quote}\small\ttfamily
19cd7f7863a681828fd026adbf8ce91c21e3ac7604439bc9d4603ba58661853e.
\end{quote}
Construct $N$ by deleting state $218$ and applying the six row changes
in Table~\ref{tab:239-edits}.  Every other row is identical.  Here
$(w,d,q)$ means write $w$, move $d$, and enter $q$.

\begin{table}[htbp]
\centering
\begin{tabular}{rrll}
State & Read & Old transition & New transition \\
\hline
$179$ & $1$ & $(1,R,105)$ & $(1,L,215)$ \\
$187$ & $1$ & $(1,R,218)$ & $(1,R,179)$ \\
$220$ & $1$ & $(1,R,218)$ & $(1,R,179)$ \\
$223$ & $0$ & $(0,R,16)$ & $(0,R,23)$ \\
$223$ & $1$ & $(1,R,99)$ & $(1,R,70)$ \\
$237$ & $0$ & $(0,R,218)$ & $(0,R,223)$ \\
\hline
\end{tabular}
\caption{All six changed rows.  The two rows of old state $218$,
$(0,L,223)$ and $(1,R,70)$, are removed.}
\label{tab:239-edits}
\end{table}

The resulting inherited-label table has $239$ working states and $478$
rows.  All $239$ working states are reachable in the directed transition
graph from state $95$.  The table file
\texttt{boundary239\_fold179\_old.tm} has SHA-256 digest
\begin{quote}\small\ttfamily
d061770f487937f9da652919c9bc9bb79bc9bf137ca41dd7a08eb5ecda8e7719.
\end{quote}
Graph reachability is used here only to account for the literal states;
it does not claim that every state is reached by the blank execution.

\subsection{The invariant and control representation}

Write $\ell=T(h-1)$ for the tape bit immediately left of the head.
Inspection of the complete old transition table establishes the
following incoming-port facts:
\begin{itemize}
\item The only incoming transitions to old $179$ are
$161/0$, $236/0$, and $239/0$.  Each writes $0$ and moves right.
Thus every entry to old $179$ has $\ell=0$.
\item The only incoming transitions to old $218$ are $187/1$,
$220/1$, and $237/0$.  The first two write $1$ and move right;
the last writes $0$ and moves right.  They therefore supply tags
$1$, $1$, and $0$, respectively.
\item The only incoming transition to old $223$ is old $218/0$.
Old $223$ can be kept entirely inside a three-step old-machine block.
\end{itemize}
None of these special states is the initial state.  The fixed tag at
old $179$ is therefore a globally inductive fact, valid even for an
arbitrary initial binary tape.  The two tags at old $218$ are exhaustive.

At simulation checkpoints the head and the entire tape are identical.
The control relation is
\[
\begin{array}{c|c}
\text{Old control and tag} & \text{New control} \\
\hline
179,\ \ell=0 & 179\\
218,\ \ell=0 & 223\\
218,\ \ell=1 & 179\\
q\notin\{179,218,223\} & q.
\end{array}
\]
Old $223$ is not a checkpoint.  Its executions occur inside the
three-step block that starts at old $218$ reading $0$.
The new targets of the three old $218$ callers select exactly the
corresponding entry of this relation: $237/0$ selects new $223$,
while $187/1$ and $220/1$ select new $179$.

\subsection{Six whole-tape block identities}

The local verification needs only the two tape cells at relative
positions $-1$ and $0$.  The cells outside this pair are arbitrary.
All six old and new blocks in Table~\ref{tab:239-blocks} leave the
\emph{entire tape unchanged}, finish with head displacement $+1$,
and enter the stated common endpoint.  No block reads or writes
outside those two cells.

\begin{table}[htbp]
\centering
\begin{tabular}{rrrrrrr}
Old state & $\ell$ & Read & New state & Old steps & New steps & Endpoint \\
\hline
$179$ & $0$ & $0$ & $179$ & $1$ & $1$ & $105$\\
$179$ & $0$ & $1$ & $179$ & $1$ & $5$ & $105$\\
$218$ & $0$ & $0$ & $223$ & $3$ & $1$ & $23$\\
$218$ & $0$ & $1$ & $223$ & $1$ & $1$ & $70$\\
$218$ & $1$ & $0$ & $179$ & $3$ & $1$ & $105$\\
$218$ & $1$ & $1$ & $179$ & $1$ & $3$ & $70$\\
\hline
\end{tabular}
\caption{The complete local proof obligations for the context split
and fold.  Each row quantifies over every surrounding binary tape.}
\label{tab:239-blocks}
\end{table}

The useful old projection is particularly simple.  At old $218$
reading $0$, the machine first moves left into old $223$.  If the
tag is $0$, it then passes through $16$ and reaches $23$.  If the
tag is $1$, it passes through $99$ and reaches $105$.  Both paths
take three steps, preserve the two bits, and finish one cell to the
right of the starting head.  New $223/0$ and new $179/0$ implement
these respective outcomes directly.

The shared new row $179/1=(1,L,215)$ uses the preceding tag to recover
the two remaining meanings.  For tag $0$ its successive executed rows are
\[
179/1,\quad 215/0,\quad 183/1,\quad 213/0,\quad 237/1,
\]
and it finishes at $105$.  For tag $1$ they are
\[
179/1,\quad 215/1,\quad 233/1,
\]
and it finishes at $70$.  The first path takes five steps and
simulates old $179/1$; the second takes three steps and simulates
old $218/1$ in its tag-$1$ context.  The tag-$0$ context of old
$218/1$ is handled directly by new $223/1$.

\subsection{Composition theorem}

\paragraph{Theorem.}
For every binary initial tape $T$, the old repaired boundary240
machine started in state $95$ on $T$ halts if and only if the
new $239$-state machine started in state $95$ on $T$ halts.
In particular, their blank-tape halting and non-halting statements
are equivalent.

\paragraph{Proof.}
The initial configurations are related because state $95$ is unchanged.
At an ordinary checkpoint, the next old and new transition have identical
writes and moves.  Their targets are either identical ordinary states,
old $179$ with its forced tag $0$, or one of the three rerouted old
$218$ ports.  The incoming-port checks above establish the control
relation in every case.  There is no ordinary transition into old $223$.

At a special checkpoint, select the corresponding row of
Table~\ref{tab:239-blocks}.  It provides a finite old block and a
finite new block ending in exactly the same ordinary control state,
head position, and tape.  Hence the checkpoint relation is restored.
This also accounts for every visit to old $223$; it cannot be entered
except through one of the specified old $218/0$ blocks.

Each block on each side has positive length.  By induction the two
deterministic executions have corresponding checkpoints for every
finite number of blocks.  A finite halting execution corresponds
to a finite halting execution, while an infinite execution supplies
infinitely many positive blocks on both sides.  The machines therefore
halt together.  All invariant checks and local identities are finite
facts about the displayed tables, and the induction is formalizable
in ordinary elementary arithmetic.  In particular, ZFC proves the
blank-tape equivalence. \hfill$\square$

Consequently, any proof in ZFC that the repaired boundary240 machine
does not halt exactly when $\operatorname{Con}(\mathrm{ZFC})$ holds
transfers to the explicit $239$-state table.  This local theorem does
not by itself discharge the separate source-calculus and compiler
correctness obligations for the repaired boundary240 reference.

\subsection{Reproducible certificate}

The standalone Python program \texttt{certify\_239\_fold.py} validates
both tables, checks the entire edit set and every incoming port of
$179$, $218$, and $223$, and executes the six blocks above.  Its
symbolic tape consists of the two specified cells; an attempted
read outside them is an error.  It checks the complete final tape
map, head, and control state, and records both step-by-step traces.
It uses only the Python standard library.  The command
\begin{verbatim}
python certify_239_fold.py boundary240_old.tm \
  boundary239_fold179_old.tm --output fold239_certificate.json
\end{verbatim}
returns \texttt{PASS}.  Alternatively, \texttt{--build} constructs the
new table directly from the old table and verifies it before saving.
The supplied \texttt{fold239\_certificate.json} records all six
obligations, the incoming-port checks, and the input hashes.  This
is a finite executable certificate together with the composition
proof above, rather than a proof-assistant kernel certificate.

The discovery process first tested $50,317$ finite runs, including
startup, collecting $26,842,232$ transition observations and all
$480$ old rows.  These observations suggested the fold.  The proof
above replaces that heuristic evidence: it uses a global entry-tag
invariant and universally quantified two-cell identities.  A separate
finite regression of all $8,386$ instruction guards on the $239$-state
table also passes.  Those finite tests are supporting diagnostics;
they are not the basis of the equivalence theorem.
